\providecommand{\Lean}{Lean~4}

\ifdefined\leandefinition\else
\newenvironment{leandefinition}[1]{%
\par\medskip\noindent\textbf{Definition (#1).}\ }{\par\medskip}
\fi
\ifdefined\leanremark\else
\newenvironment{leanremark}[1]{%
\par\medskip\noindent\textbf{Remark (#1).}\ }{\par\medskip}
\fi

\section{Lean}
\label{sec:lean}

\Lean{} is both a dependently typed functional programming language and an
interactive theorem prover (ITP)~\cite{deMouraUllrich2021Lean4,LeanReferenceManual}.  The two aspects are not independent layers:
programs, specifications, and proofs are expressions of the same typed language.
A theorem is represented by a constant whose type is a proposition, and a proof
is represented by a term inhabiting that type.  Tactics, automation, notation,
and elaboration are user-facing mechanisms for constructing such terms, but the
trusted kernel only checks fully elaborated proof terms.

Lean contains automated procedures, and these procedures may solve many goals without further user input.
Nevertheless, the system is interactive because the user chooses statements, definitions, intermediate lemmas, automation strategies, and proof
structure.
Moreover, automation is not accepted as an oracle: every tactic must produce a term that is checked by the kernel.

\subsection{Logical foundations}
\label{subsec:lean-foundations}

Lean's kernel is based on an intensional dependent type theory in the family of
the Calculus of Inductive Constructions (CIC)~\cite{coquandHuet1988coc,coquandPaulin1990cic}.
We recall the features that the developments of this thesis rely on.

\paragraph{Universes.}
Because every type is itself a term with a type of its own, the types cannot all
inhabit a single type: the rule \leaninl{Type : Type} would render the system
inconsistent, by Girard's paradox~\cite{girard1972}---a type-theoretic analog of
Russell's paradox in set theory (as explained in \Cref{subsec:set-theory-history}).
Lean therefore stratifies types into a non-cumulative hierarchy of \emph{universes}
\leaninl{Sort u}, indexed by a \emph{universe level} \leaninl{u}, each inhabiting
the next: \leaninl{Sort u : Sort (u+1)}.

Two named families sit inside this one hierarchy. At the bottom is the universe of
propositions \leaninl{Prop}, defined as \leaninl{Sort 0}; above it are the
\emph{data} universes \leaninl{Type u}, abbreviating \leaninl{Sort (u+1)}. Thus,
\leaninl{Type} and \leaninl{Sort} name the same tower shifted by one index: small
types such as \leaninl{Nat : Type} and \leaninl{Prop : Type} sit low in it,
whereas a constructor quantifying over a universe lands one level higher. Unlike set theory,
the universe hierarchy is \emph{non-cumulative}---a type in \leaninl{Type u} does not silently
belong to \leaninl{Type (u+1)}---so movement between levels is explicit.

\paragraph{Universe polymorphism.}
A construction that makes sense at every level is declared once with a universe
\emph{variable} \leaninl{u} rather than duplicated level by level. This is the
role of the \leaninl{u} appearing in the examples below, such as the type
\leaninl{Tree} of \Cref{ex:lean-inductive}, whose \leaninl{Type u} ranges over all
universes.

\paragraph{Propositions as types.}
\leaninl{Prop} differs from the data universes in two respects that reflect its
logical role. It is \emph{impredicative}---a quantified proposition
\leaninl{∀ x : A, p x} stays in \leaninl{Prop} regardless of the universe of
\leaninl{A}---and the kernel treats it as \emph{proof-irrelevant}: any two proofs
of one proposition are interchangeable, since only provability, not the identity
of the proof, carries information. Under the Curry--Howard correspondence, a
proposition \leaninl{p : Prop} is the type of its proofs, a proof is a term
\leaninl{h : p} inhabiting that type, and the logical connectives are ordinary inductive
types---conjunction a product, disjunction a sum, implication a function
type. Proving a theorem and constructing a term of a given type are therefore the
same activity, carried out either by writing the term or by running tactics that
elaborate one.

\paragraph{Inductive types.}
An inductive type is specified by its constructors, from which the kernel
derives the matching recursion and induction principles. Data types and logical
connectives are alike instances of this single mechanism---the reason programs
and proofs inhabit one language.

\paragraph{Large elimination.}
\label{par:large-elimination}
Every inductive type \leaninl{T} comes with a \emph{recursor} \leaninl|T.rec|
universe-polymorphic in the level \leaninl{v} of the result it produces (the elided
hypotheses are the minor premises, one per constructor):
\begin{leanlisting}
T.rec.{v} : (motive : T → Sort v) → ... → (t : T) → motive t
\end{leanlisting}
When \leaninl{T} is a data type, \leaninl{v} is unconstrained.
When \leaninl{T : Prop}, the kernel instead
forces \leaninl{v = 0}: a proof may be eliminated only into another proposition, a
restriction called \emph{small elimination}. Eliminating a proposition into a data
universe is \emph{large elimination}, and it is forbidden in general because it is
inconsistent with the proof irrelevance of \leaninl{Prop}---from a single proof one
could compute which constructor produced it, telling apart proofs the kernel deems
equal. The one exception is a \emph{subsingleton}---an inductive with at most one
constructor, all of whose fields are themselves propositions (\leaninl{Eq},
\leaninl{And}, \leaninl{PUnit}, \etc)---whose recursor may target any universe.

\begin{example}[a forbidden universe-polymorphic inductive]
\label{ex:large-elimination}
The following declaration is rejected by the kernel:
\begin{leanlisting}
inductive T : Sort u
\end{leanlisting}
Indeed, at \leaninl{u := 0} the type \leaninl{T} would be a proposition,
whose recursor would then large-eliminate---unsoundly. Lean therefore
refuses the polymorphic \leaninl{Sort u} and requires either \leaninl{Prop}
or a data universe \leaninl{Type u}.
\end{example}

\paragraph{Classical and computable fragments.}
Lean is constructive by default but \emph{classical in the presence of choice}:
by Diaconescu's theorem~\cite{diaconescu1975}, the choice axiom together with
functional extensionality entails the law of excluded middle. A development may
nonetheless remain within a \emph{computable} fragment by avoiding choice and the
\leaninl{Classical} module; a definition that does depend on choice cannot be
reduced by the kernel, which then requires it to be marked \leaninl{noncomputable}.
The operator used in later chapters is \leaninl{Classical.choose}, which turns a proof of
\leaninl{∃ x, p x} into a witness, together with \leaninl{Classical.choose_spec}, which
states that the witness satisfies \leaninl{p}.

\subsection{Lean as a functional programming language}
\label{subsec:lean-programming}

Lean is a pure, strict, functional language with dependent types.
The logical fragment is total: every function that is used as a definitional object in the logic must be justified as terminating, typically by structural recursion or by a well-founded recursion argument.
Lean still provides pragmatic mechanisms such as \leaninl{partial} and \leaninl{unsafe}, but these mechanisms are separated from the trusted logical fragment.

Lean as a programming language is convenient for formalization for three main reasons.
Firstly, executable algorithms and their specifications are written in the same
language.
Secondly, dependent types allow ordinary programming interfaces to
include logical invariants, for example the bound of an array index.
Thirdly, the meta-programming language is Lean itself: tactics, commands, syntax
extensions, and domain-specific proof procedures can be implemented in the same
environment as the formal development.

\begin{example}[inductive data and recursion]
\label{ex:lean-inductive}
Defining a polymorphic type for representing binary trees and a recursive function counting their internal nodes can be done as follows.

\begin{leanlisting}
inductive Tree (α : Type u) : Type u where
  | leaf : Tree α
  | node : Tree α → α → Tree α → Tree α
\end{leanlisting}

When Lean processes an inductive type like \leaninl{Tree}, it also generates a case elimination and induction principles---for both proofs and programs.

\begin{leanlisting}
def Tree.size {α : Type u} : Tree α → Nat
  | .leaf => 0
  | .node l _ r => Tree.size l + 1 + Tree.size r
\end{leanlisting}
The recursive calls are made on strict subtrees; hence the definition is structurally recursive.
\end{example}

Dependent types are useful even in small programs.
\Cref{ex:array-bounds} illustrates a common Lean pattern: some types may carry proof components that are relevant for type checking, but since they are propositions, they are erased from compiled code.

\begin{example}
\label{ex:array-bounds}
If \leaninl{xs : Array α},
then \leaninl{Fin xs.size} is the type of valid indices into \leaninl{xs}.  A
function taking such an index does not need a dynamic bounds check in its
specification: the proof of boundedness is part of the input type.

\begin{leanlisting}
def indicesOf {α : Type u} (xs : Array α) (x : α) : Set α :=
  { i : Fin xs.size | xs[i] = x }
\end{leanlisting}

If validity of an index cannot be determined by the kernel automatically, the proof of its validity can still be supplied via \leaninl{xs[i]'hi}, where \leaninl{hi : i < xs.size} is a proof of the index's validity.

\end{example}

\subsection{Lean as a proof assistant}
\label{subsec:lean-proof-assistant}

Lean implements the Curry--Howard correspondence: propositions are types and
proofs are terms.  If \leaninl{p : Prop}, then \leaninl{p} is itself a type, and a proof
of \leaninl{p} is a term \leaninl{h : p} inhabiting it.  A theorem declaration is therefore
a definition whose result type is
in \leaninl{Prop}.  Lean supports both term-style proofs, which explicitly build proof
terms, and tactic-style proofs, where commands incrementally transform goals and
construct a term implicitly.

\begin{example}[term-style and tactic-style proofs]
Both theorems below state the same property, whose proofs are written in term style and tactic style respectively.

\begin{leanlisting}
theorem and_swap (p q : Prop) : p ∧ q → q ∧ p :=
  fun h ↦ And.intro h.right h.left

theorem and_swap' (p q : Prop) : p ∧ q → q ∧ p := by
  intro h
  apply And.intro
  · exact h.left
  · exact h.right
\end{leanlisting}
In both cases the proposition \leaninl{p ∧ q → q ∧ p} is a type living in \leaninl{Prop}, and the term-style \leaninl{and_swap} is literally a function: it takes an argument \leaninl{h : p ∧ q} and returns a term of type \leaninl{q ∧ p}.
\end{example}
